\documentclass[11pt]{article}
\usepackage[a4paper,margin=2.2cm]{geometry}
\usepackage{xcolor}
\usepackage{fontspec}
\usepackage{xeCJK}
\usepackage{enumitem}
\setmainfont{Times New Roman}
\setCJKmainfont{SimSun}
\setlist[itemize]{leftmargin=1.4em,itemsep=0.35em,topsep=0.35em}
\setlength{\parindent}{0pt}
\setlength{\parskip}{0.55em}
\pagestyle{plain}
\begin{document}
{\color[HTML]{1F5FD2}\LARGE\bfseries 特殊教育评估工具使用手册\par}
{\color[HTML]{667085}\small 星轨原创教育资源：用于家庭与学校共同制定支持计划。\par}
\vspace{0.8em}
{\color[HTML]{111827}\large\bfseries 评估目的\par}
\begin{itemize}
\item 特殊教育评估用于理解儿童当前能力、参与障碍和支持需求，目的是指导教学，而不是降低期待。
\item 评估结论应能回答下一步教什么、怎样教、由谁支持、多久复盘。
\end{itemize}
{\color[HTML]{111827}\large\bfseries 多元资料\par}
\begin{itemize}
\item 可结合标准化测验、课程本位测量、课堂观察、家庭访谈和作品样本，避免只依赖单一工具。
\item 当语言、文化、焦虑或动作操作影响表现时，更需要在不同情境中验证结果。
\end{itemize}
{\color[HTML]{111827}\large\bfseries 结果转化\par}
\begin{itemize}
\item 评估发现要转化为可教学的支持。例如听觉工作记忆弱，可转化为书面步骤、复述检查和视觉提醒。
\item 报告建议应写成可执行行动，明确负责人、频率和观察指标。
\end{itemize}
\vfill
{\color[HTML]{667085}\footnotesize 本材料为应用内原创教育内容，不能替代医疗、法律或正式评估建议。\par}
\end{document}
